\documentclass[10pt,a4paper]{amsart}
\usepackage{amsmath,amssymb,amsthm,graphicx}
\usepackage{verbatim}
\newenvironment{reference}{\comment}{\endcomment}
%\newenvironment{reference}{}{}
\newenvironment{slogan}{\comment}{\endcomment}
\newenvironment{history}{\comment}{\endcomment}

% For commutative diagrams we use Xy-pic
\usepackage[all]{xy}

% We use 2cell for 2-commutative diagrams.
\xyoption{2cell}
% Theorem environments.
%
\theoremstyle{plain}
\newtheorem{theorem}[subsection]{Theorem}
\newtheorem{proposition}[subsection]{Proposition}
\newtheorem{lemma}[subsection]{Lemma}

\theoremstyle{definition}
\newtheorem{definition}[subsection]{Definition}
\newtheorem{example}[subsection]{Example}
\newtheorem{exercise}[subsection]{Exercise}
\newtheorem{situation}[subsection]{Situation}

\theoremstyle{remark}
\newtheorem{remark}[subsection]{Remark}
\newtheorem{remarks}[subsection]{Remarks}

\numberwithin{equation}{subsection}

% Macros
%
\def\lim{\mathop{\mathrm{lim}}\nolimits}
\def\colim{\mathop{\mathrm{colim}}\nolimits}
\def\Spec{\mathop{\mathrm{Spec}}}
\def\Hom{\mathop{\mathrm{Hom}}\nolimits}
\def\Ext{\mathop{\mathrm{Ext}}\nolimits}
\def\SheafHom{\mathop{\mathcal{H}\!\mathit{om}}\nolimits}
\def\SheafExt{\mathop{\mathcal{E}\!\mathit{xt}}\nolimits}
\def\Sch{\mathit{Sch}}
\def\Mor{\mathop{\mathrm{Mor}}\nolimits}
\def\Ob{\mathop{\mathrm{Ob}}\nolimits}
\def\Sh{\mathop{\mathit{Sh}}\nolimits}
\def\NL{\mathop{N\!L}\nolimits}
\def\CH{\mathop{\mathrm{CH}}\nolimits}
\def\proetale{{pro\text{-}\acute{e}tale}}
\def\etale{{\acute{e}tale}}
\def\QCoh{\mathit{QCoh}}
\def\Ker{\mathop{\mathrm{Ker}}}
\def\Im{\mathop{\mathrm{Im}}}
\def\Coker{\mathop{\mathrm{Coker}}}
\def\Coim{\mathop{\mathrm{Coim}}}

% Boxtimes
%
\DeclareMathSymbol{\boxtimes}{\mathbin}{AMSa}{"02}

%
% Macros for moduli stacks/spaces
%
\def\QCohstack{\mathcal{QC}\!\mathit{oh}}
\def\Cohstack{\mathcal{C}\!\mathit{oh}}
\def\Spacesstack{\mathcal{S}\!\mathit{paces}}
\def\Quotfunctor{\mathrm{Quot}}
\def\Hilbfunctor{\mathrm{Hilb}}
\def\Curvesstack{\mathcal{C}\!\mathit{urves}}
\def\Polarizedstack{\mathcal{P}\!\mathit{olarized}}
\def\Complexesstack{\mathcal{C}\!\mathit{omplexes}}
% \Pic is the operator that assigns to X its picard group, usage \Pic(X)
% \Picardstack_{X/B} denotes the Picard stack of X over B
% \Picardfunctor_{X/B} denotes the Picard functor of X over B
\def\Pic{\mathop{\mathrm{Pic}}\nolimits}
\def\Picardstack{\mathcal{P}\!\mathit{ic}}
\def\Picardfunctor{\mathrm{Pic}}
\def\Deformationcategory{\mathcal{D}\!\mathit{ef}}
\usepackage{cleveref}
\crefname{originallemma}{original lemma}{original lemmas}
\Crefname{originallemma}{Original lemma}{Original lemmas}
\makeatletter
\newlabel{algebra-lemma-hathat-finitely-generated}{{05GG}{source}{}{sourcepointer-1}{}}
\newlabel{algebra-lemma-hathat-finitely-generated@cref}{{[originallemma][0][]05GG}{[1][0][]source}{}{}{}}
\makeatother
\crefname{originallemma}{original lemma}{original lemmas}
\Crefname{originallemma}{Original lemma}{Original lemmas}
\makeatletter
\newlabel{algebra-lemma-quotient-complete}{{031A}{source}{}{sourcepointer-2}{}}
\newlabel{algebra-lemma-quotient-complete@cref}{{[originallemma][0][]031A}{[1][0][]source}{}{}{}}
\makeatother
\title{A ring complete for a principal ideal with a principal ideal whose quotient is not complete}
\author{The Stacks Project Authors}
\date{stacks-project:examples.tex:local-snapshot-20261012; GNU FDL 1.2 or later}
\begin{document}
\maketitle
\noindent\textit{Editorial note.} All delivered mathematical text is the author's own, copied from the pinned Stacks Project chapter file examples.tex; no statement or proof step was rewritten, reordered or removed.\par\medskip
\noindent\textit{Editorial note.} The counted proof is the author's complete argument: The counted proof adds no mathematics of its own; the complete chain (canonical form, replacement rules, contradiction) lives in the context span.\par\medskip
\noindent\textit{Editorial note.} The label lemma-noncomplete-quotient names the counted statement, and the proof environment immediately following it in the source is delivered verbatim as its proof.\par\medskip
\noindent\textit{Editorial note.} The counted statement is carried by the author's own lemma environment in the chapter source.\par\medskip
\noindent\textit{Allowed standard prerequisites (editorial).}\begin{itemize}
\item Context span 680-727: R = k[t, z\_1, z\_2, ..., w\_1, w\_2, ..., x]/(z\_i t - x\textasciicircum{}i w\_i, z\_i w\_j); every element has a unique canonical form f = sum\_\{i=0..d\} f\_i x\textasciicircum{}i with f\_i in k[t, z\_i, w\_j] avoiding the products z\_i t and z\_i w\_j; f in (x\textasciicircum{}n) iff f\_i = 0 for i < n; so x is a nonzerodivisor and the intersection of the (x\textasciicircum{}n) is 0
\item The (x)-adic completion R\textasciicircum{}wedge consists of unique infinite sums f = sum\_\{i>=0\} f\_i x\textasciicircum{}i (Algebra, hathat-finitely-generated); the element f = sum\_\{i>=1\} x\textasciicircum{}i w\_i lies in (t, x\textasciicircum{}n) for every n because x\textasciicircum{}m w\_m is in (t); the claim f not in (t) is proven by the replacement rules for multiplying t g in canonical form (a term of g\_0 involving z\_i contributes c w\_i to f\_\{l+i\}; picking n with z\_n absent from the polynomial g\_0 leaves w\_n absent from f\_n, a contradiction)
\item Algebra, quotient-complete: a quotient of a complete ring is complete iff it is classically complete, so R\textasciicircum{}wedge/(t) is not complete
\end{itemize}\medskip
\ifdefined\hypertarget\hypertarget{sourcepointer-1}{}\fi
\noindent\textit{Source reference.} Original lemma 05GG: sources/stacks-chow/tags.txt, source lines 7060--7060. Prerequisite pointer invoked at 755; it names a standard fact of a sibling chapter that the counted argument uses, and every step of the counted proof built on it is retained verbatim in the delivered body.\par\medskip
\ifdefined\hypertarget\hypertarget{sourcepointer-2}{}\fi
\noindent\textit{Source reference.} Original lemma 031A: sources/stacks-chow/tags.txt, source lines 3977--3977. Prerequisite pointer invoked at 786; it names a standard fact of a sibling chapter that the counted argument uses, and every step of the counted proof built on it is retained verbatim in the delivered body.\par\medskip
Let $k$ be a field. Let
$$
R = k[t, z_1, z_2, z_3, \ldots, w_1, w_2, w_3, \ldots, x]/
(z_it - x^iw_i, z_i w_j)
$$
Note that in particular $z_iz_jt = 0$ in this ring. Any element $f$ of $R$
can be uniquely written as a finite sum
$$
f = \sum\nolimits_{i = 0, \ldots, d} f_i x^i
$$
where each $f_i \in k[t, z_i, w_j]$ has no terms involving the products
$z_it$ or $z_iw_j$. Moreover, if $f$ is written in this way, then
$f \in (x^n)$ if and only if $f_i = 0$ for $i < n$.
So $x$ is a nonzerodivisor and $\bigcap (x^n) = 0$.
Let $R^\wedge$ be the completion of $R$ with respect to the ideal $(x)$.
Note that $R^\wedge$ is $(x)$-adically complete, see
Algebra, Lemma \ref{algebra-lemma-hathat-finitely-generated}.
By the above we see that an element of $R^\wedge$ can be uniquely written
as an infinite sum
$$
f = \sum\nolimits_{i = 0}^\infty f_i x^i
$$
where each $f_i \in k[t, z_i, w_j]$ has no terms involving the products
$z_it$ or $z_iw_j$. Consider the element
$$
f = \sum\nolimits_{i = 1}^\infty x^i w_i =
xw_1 + x^2w_2 + x^3w_3 + \ldots
$$
i.e., we have $f_n = w_n$. Note that $f \in (t , x^n)$ for every $n$
because $x^mw_m \in (t)$ for all $m$.
We claim that $f \not \in (t)$. To prove this assume that
$tg = f$ where $g = \sum g_lx^l$ in canonical form as above.
Since $tz_iz_j = 0$ we may as well assume that none of the $g_l$ have
terms involving the products $z_iz_j$. Examining the process to
get $tg$ in canonical form we see the following:
Given any term $c m$ of $g_l$ where $c \in k$ and $m$ is a
monomial in $t, z_i, w_j$ and we make the following replacement
\begin{enumerate}
\item if the monomial $m$ does not involve any $z_i$, then $ctm$ is
a term of $f_l$, and
\item if the monomial $m$ does involve a $z_i$ then it is equal to
$m = z_i$ and we see that $cw_i$ is term of $f_{l + i}$.
\end{enumerate}
Since $g_0$ is a polynomial only finitely many of the variables $z_i$
occur in it. Pick $n$ such that $z_n$ does not occur in $g_0$.
Then the rules above show that $w_n$ does not occur in $f_n$ which is
a contradiction. It follows that $R^\wedge/(t)$ is not complete, see
Algebra, Lemma \ref{algebra-lemma-quotient-complete}.

\begin{lemma}
\label{lemma-noncomplete-quotient}
There exists a ring $R$ complete with respect to a principal ideal
$I$ and a principal ideal $J$ such that $R/J$ is not $I$-adically
complete.
\end{lemma}

\begin{proof}
See discussion above.
\end{proof}

\end{document}
